\documentclass[12pt,a4paper]{article}
\usepackage[utf8]{inputenc}
\usepackage[T1]{fontenc}
\usepackage[margin=2.5cm]{geometry}
\usepackage{graphicx}
\usepackage{enumitem}
\usepackage{hyperref}
\usepackage{xcolor}
\usepackage{tcolorbox}

\title{\textbf{Getting Google API Keys for the App}}
\author{}
\date{}

\setlist{itemsep=4pt}
\hypersetup{colorlinks=true, urlcolor=blue}

\begin{document}

\maketitle
\vspace{-2em}

\section*{Before you start -- two-factor authorization}

When you go to \url{https://console.cloud.google.com} for the first time, it
will ask that your email account be set up with \textbf{two-factor
authorization}. The easiest way to do this is by installing the
\textbf{Google Authenticator} app on your phone.

\begin{enumerate}
  \item Install \textbf{Google Authenticator} from your phone's app store
        (Google Play / App Store).
  \item Log in with the \textbf{same email} in both Google Authenticator and
        the website.
  \item Verify Google Authenticator on your phone.
  \item After that, \textbf{enable two-factor authorization} on your account.
  \item Once enabled, you will be able to access the website.
\end{enumerate}

After you can access the website, follow the steps below.

\section*{What you need to do}

To get the map working in the app, we need \textbf{three API keys} from
Google (plus one optional key for a website). Please follow the steps below
and send the keys to the developer. Each key is a long string that starts
with \texttt{AIza}.

\begin{enumerate}
  \item Go to \url{https://console.cloud.google.com} and sign in with your
        Google account.
  \item Create a new project (or choose an existing one). Click the project
        name at the top of the page, then \textbf{New Project}, give it a
        name, and click \textbf{Create}.
  \item After the project is created, make sure it is selected (you should see
        its name at the top of the screen).
  \item Go to \url{https://console.cloud.google.com/billing} and \textbf{enable
        billing} on the project. \textbf{Don't worry} -- a map-only app costs
        \$0. Billing just needs to be enabled even though nothing will be
        charged.
\end{enumerate}

\includegraphics[width=0.95\textwidth]{1.jpg}
\begin{center}
\textit{Image 1: Project and billing setup}
\end{center}

\section*{Key 1 -- Map SDK for Android}

\begin{enumerate}
  \item Open \url{https://console.cloud.google.com/apis/library/maps-android-backend.googleapis.com}
        and click \textbf{Enable}.
  \item Go to \url{https://console.cloud.google.com/apis/credentials} and click
        \textbf{Create Credentials} $>$ \textbf{API Key}.
  \item Copy the key that appears (it starts with \texttt{AIza}).
  \item Label it \textbf{Android} so you can tell them apart.
\end{enumerate}

\includegraphics[width=0.95\textwidth]{2.jpg}
\begin{center}
\textit{Image 2: Creating an API key}
\end{center}

\section*{Key 2 -- Map SDK for iOS}

\begin{enumerate}
  \item Open \url{https://console.cloud.google.com/apis/library/maps-backend.googleapis.com}
        and click \textbf{Enable}.
  \item Go to \url{https://console.cloud.google.com/apis/credentials} and click
        \textbf{Create Credentials} $>$ \textbf{API Key}.
  \item Copy the key (starts with \texttt{AIza}).
  \item Label it \textbf{iOS}.
\end{enumerate}

\section*{Key 3 -- Map SDK for Directions}

\begin{enumerate}
  \item Open \url{https://console.cloud.google.com/apis/library/directions-backend.googleapis.com}
        and click \textbf{Enable}.
  \item Go to \url{https://console.cloud.google.com/apis/credentials} and click
        \textbf{Create Credentials} $>$ \textbf{API Key}.
  \item Copy the key (starts with \texttt{AIza}).
  \item Label it \textbf{Directions}.
\end{enumerate}

\section*{Key 4 -- Maps JavaScript API \textit{(optional, for a website)}}

This key is \textbf{not needed for the phone app}. It is only for a website
later on. Since it is free with no limit, you can create it now if you want.
It is completely optional.

\begin{enumerate}
  \item Open \url{https://console.cloud.google.com/apis/library/maps-backend.googleapis.com}
        and click \textbf{Enable}.
  \item Go to \url{https://console.cloud.google.com/apis/credentials} and click
        \textbf{Create Credentials} $>$ \textbf{API Key}.
  \item Copy the key (starts with \texttt{AIza}).
  \item Label it \textbf{JavaScript}.
\end{enumerate}

\includegraphics[width=0.95\textwidth]{3.jpg}
\begin{center}
\textit{Image 3: Copying the API key}
\end{center}

\section*{Sending the keys to us}

\begin{tcolorbox}[colback=blue!5, colframe=blue!60, title=What to send]
Please send us the \textbf{three required keys}, clearly labelled (the
JavaScript key is optional):
\begin{itemize}
  \item \textbf{Android key:} \texttt{AIza...}
  \item \textbf{iOS key:} \texttt{AIza...}
  \item \textbf{Directions key:} \texttt{AIza...}
  \item \textbf{JavaScript key (optional):} \texttt{AIza...}
\end{itemize}
\end{tcolorbox}

\begin{itemize}
  \item Keep the keys private and only share them with us.
  \item The Directions API has a limited free allowance each month. If you
        expect heavy usage, let the developer know so we can plan for it.
\end{itemize}

\end{document}
